% ==============================================================================
\section{Problem Statement}
\label{sec:problem}
% ==============================================================================

We model a hangar as a directed graph $G = (\calP, \calA)$, where
$\calP = \{p_1, \dots, p_P\}$ is the finite set of \emph{maintenance positions}
and $\calA \subseteq \calP \times \calP$ is a set of \emph{blocking arcs}.  An
arc $(p, p') \in \calA$ means that position $p$ (\emph{front}) physically
blocks access to position $p'$ (\emph{rear}): an aircraft at position $p'$
cannot enter or leave the hangar without temporarily displacing the aircraft
at position $p$.  In practice $G$ is acyclic, so a position cannot block
itself transitively, and we define the \emph{blocking load} of a position as
$\ell(p) = |\{q \in \calP : p \text{ can reach } q \text{ in } G\}|$.  A
higher blocking load means that aircraft at $p$ can obstruct more other
positions, making $p$ topologically \emph{central}.

\begin{definition}[Blocking arc]
  An arc $(p, p') \in \calA$ induces an \emph{entry movement} whenever an
  aircraft arrives at position $p'$ while another aircraft is being serviced
  at position $p$, and an \emph{exit movement} whenever an aircraft departs
  from position $p'$ while another aircraft is still at position $p$.  Each
  movement contributes 2 to the movement count, one for the displacement of
  the front aircraft and one for its return.
\end{definition}

Let $\calR = \{r_1, \dots, r_R\}$ be the set of aircraft.  Each aircraft
$r \in \calR$ requires a fixed amount of maintenance work of total
duration $D_r > 0$, which we treat as atomic: while $r$ occupies a
position, the position is busy for exactly $D_r$ time units in a single
uninterrupted block.\footnote{In practice the maintenance work consists of
an ordered sequence of jobs, and the per-job timings can be reconstructed
deterministically once the aircraft start time is known
(see \S\ref{sec:milp}).  Because the operationally relevant quantities
(makespan, delay, blocking interactions) depend only on the aggregated
aircraft block, the problem can be stated entirely at the aircraft level.}
Each aircraft has an \emph{earliest start} $E_r \ge 0$ before which work
cannot begin and a \emph{target finish} $L_r > E_r$ beyond which
completion incurs delay.

A solution to the AHPP makes two coupled decisions.  At the assignment
level, it sends each aircraft $r$ to exactly one position
$\pi(r) \in \calP$.  At the scheduling level, it fixes a start time
$s_r \ge 0$ for every aircraft, from which the finish time follows as
$f_r = s_r + D_r$.

Any solution must satisfy the following requirements.  Every aircraft is
\emph{assigned to exactly one position}, and a position handles its
aircraft \emph{sequentially}: when two aircraft $r$ and $r'$ are assigned
to the same position, their occupancy intervals $[s_r, f_r]$ and
$[s_{r'}, f_{r'}]$ must not overlap and a minimum separation
$\varepsilon > 0$ must elapse between the finish of one and the start of
the next.  Each aircraft's start time must respect its
\emph{earliest-start time window}, $s_r \ge E_r$.

The blocking-arc structure imposes the final family of requirements.  For
every arc $(p, p') \in \calA$ and every pair of aircraft $r, r'$ with
$\pi(r) = p$ and $\pi(r') = p'$, an entry movement is incurred if $r'$
arrives at the rear position while $r$ is still being serviced at the
front, that is, if $s_r \le s_{r'} \le f_r - \varepsilon$; analogously,
an exit movement is incurred whenever $r'$ departs from the rear while
$r$ is still at the front, that is, if $s_r \le f_{r'} \le f_r -
\varepsilon$.  Each such event triggers a temporary relocation of the
front aircraft and is counted as two movements (displacement plus return).

\paragraph{Single $\varepsilon$ for two roles.}
The parameter $\varepsilon$ plays the same physical role in both of the
requirements above: it is the time needed to tow an aircraft in or out of
a hangar slot.  This primitive operation has the same duration whether it
is performed as a turnaround between two consecutive maintenance slots at
one position --- the role of $\varepsilon$ in the separation requirement
--- or as a temporary displacement of the front aircraft to free a
blocked rear position --- the role of $\varepsilon$ in the
$f_r - \varepsilon$ margin of the blocking conditions.  We therefore
treat the two appearances as a single parameter.  Throughout this work
all temporal magnitudes are expressed in days; for the reported benchmark
we use $\varepsilon = 0.5$ days, an indicative half-day allowance for
this primitive towing operation.  The value is recorded in each instance
file, so that experiments on hangars with different operational rhythms
can use a different $\varepsilon$ without changing the model.

The objective is a weighted combination of three criteria,
\begin{equation}
  \label{eq:obj}
  \min \;\; W^M \cdot m
        \;+\; W^D \cdot \sum_{r \in \calR} v^D_r
        \;+\; W^S \cdot n,
\end{equation}
where $m = \max_{r \in \calR} f_r$ is the makespan,
$v^D_r = \max(0,\, f_r - L_r)$ is the delay of aircraft $r$, $n$ counts
the blocking movements induced above, and $W^M, W^D, W^S \ge 0$ are
user-defined weights.

\begin{remark}
  Throughout the experiments we adopt as default
  $(W^M, W^D, W^S) = (0.1, 1.0, 10)$, which encodes the operational
  priority order \emph{movements $\succ$ delay $\succ$ makespan}.
  Section~\ref{sec:experiments} also evaluates a delay-priority profile
  $(1, 10, 0.1)$ and a makespan-priority profile $(10, 0.1, 1)$ to expose
  the sensitivity of the solutions to this choice.
\end{remark}

The AHPP is $\mathcal{NP}$-hard, since it subsumes the single-machine
weighted-completion-time scheduling problem with arbitrary release dates,
which is already $\mathcal{NP}$-hard~\citep{Lenstra1977}; the additional
position-assignment layer and the blocking-arc constraints further enlarge
the search space.  For a benchmark instance with $P = 5$ positions and
$R = 10$ aircraft, the MILP formulation of \S\ref{sec:milp} involves
several thousand binary variables and an order of magnitude more
constraints; at $R = 30$ these grow by an additional order of magnitude,
motivating the heuristic approaches of \S\ref{sec:approaches}.
